% ===========================================================================
\section{Geometry in Low Dimensions}\label{sec:lowdim}
\warmth{0}\quad\whisper{Where equations become manifolds: circles, Riemann spheres, ovals, and a torus.}

\begin{strand}
Smooth projective varieties in low dimensions \emph{are} familiar manifolds. Over $\R$ a smooth
projective variety of dimension $d$ is a compact real manifold of dimension $d$; over $\C$, of
dimension $2d$. So $\PP^1_{\R}$ is a \fillin[2cm] and $\PP^1_{\C}$ is the \keyword{Riemann sphere}. This
section reads the topology straight off the algebra (degree, genus).
\end{strand}

\block{Curves in the plane: ovals, pseudolines, genus}

\begin{yourturn}
($n=1$, binary forms.) A proper subvariety of $\PP^1_{\C}$ is the zero set of a binary form
$f(x,y)$ (homogeneous in two variables). For $f=x^{11}y-11x^6y^6-xy^{11}$, the variety $\V(f)$ has
$\dim 0$ and degree $\fillin[1cm]$. These $12$ points on the Riemann sphere are Klein's
\keyword{icosahedron} vertices: $4$ real, $8$ in conjugate pairs. \recall{Degree of a binary form $=$ its number of roots in $\PP^1_{\C}$, with multiplicity.}
\end{yourturn}

For a curve $C\subset\PP^2_{\R}$, a connected component is a \keyword{pseudoline} if removing it leaves
$\PP^2_{\R}$ connected, and an \keyword{oval} otherwise (it has an inside and an outside, like a conic).

\begin{theorem}[curves in the plane, Thm.\ 2.26]\label{thm:curve}
Let $C$ be a smooth curve of degree $d$ in $\PP^2_K$.
\begin{itemize}[leftmargin=1.3em,itemsep=1pt]
  \item If $K=\C$: $C$ is an orientable surface of \keyword{genus} $g=\dfrac{(d-1)(d-2)}{2}$.
  \item If $K=\R$: $C$ has at most $g+1$ components; if $d$ is even all are ovals, if $d$ is odd exactly
        one is a \fillin[2.5cm].
\end{itemize}
\end{theorem}
\begin{yourturn}
Fill the genus by degree: $d=1\Rightarrow g=\fillin[0.6cm]$, $\ d=2\Rightarrow g=\fillin[0.6cm]$,
$\ d=3\Rightarrow g=\fillin[0.6cm]$. So lines and conics are spheres ($g=0$), and a smooth cubic is a
\fillin[1.5cm] ($g=1$).
\end{yourturn}

\block{Elliptic curves: a cubic that is a torus}

\begin{strand}
A smooth cubic ($d=3,g=1$), e.g.\ the Weierstrass curve $zy^2=x^3-xz^2$, is an \keyword{elliptic
curve}. Over $\R$ it shows an oval and a pseudoline; over $\C$ it is homeomorphic to a
\keyword{torus} $S^1\times S^1=\R^2/M$, where $M=\Z\lambda+\Z\tau$ is the period lattice obtained by
integrating $\int \mathrm dx/y$ along loops. The torus carries a group law, the seed of $19$th-century
geometry and modern elliptic-curve cryptography. \recall{Beware two ``tori'': this \emph{topological} torus is \emph{not} the algebraic torus $(\C^\ast)^n$ of toric geometry (Ch.\,8). An elliptic curve is not toric.}
\end{strand}

\block{The dictionary, completed}

\begin{strand}
Look back at the front-page table and check every row you can now fill. The two chapters were one
conversation, in two languages:
\end{strand}
\noindent\begin{tabularx}{\linewidth}{@{}L L@{}}
\toprule
{\color{indigo}Algebra} & {\color{madder}Geometry}\\
\midrule
prime ideal & irreducible variety \;(Prop.\ \ref{prop:irr-prime})\\
maximal ideal & a point\\
$I\subseteq J$ & $\V(I)\supseteq\V(J)$ \;(inclusion-reversing)\\
$I+J$ \;/\; $I\cap J$ & intersection \;/\; union\\
dimension, degree of $I$ & dimension, degree of $\V(I)$ \;(they agree!)\\
homogeneous ideal & projective variety\\
$\Ideal(\V(I))=\rad(I)$ & the Nullstellensatz bridge \;(Ch.\,6)\\
primary decomposition & decomposition into irreducibles \;(Ch.\,3)\\
\bottomrule
\end{tabularx}

\begin{strand}
Where this goes next. Chapter 3 (\emph{Solving and Decomposing}) turns ``$\V(I)=$ union of
irreducibles'' into an algorithm: \keyword{primary decomposition}, the geometry behind the Chinese
Remainder Theorem you already met in the running cubic. Keep the cubic, the dictionary, and the
Gröbner engine in hand; everything later is built from these three.
\end{strand}

\loose{Exercise 15 (book): solve Diophantus' problem -- find a rational point on the elliptic curve $y(6-y)=x^3-x$. Hint: take the tangent line at $(-1,0)$ and reintersect.}
